\documentclass[UTF8,a4paper,11pt]{ctexart}
\usepackage[margin=2cm]{geometry}
\usepackage{amsmath,amssymb,fancyhdr}
\pagestyle{fancy}\fancyhf{}\fancyfoot[C]{\thepage}\renewcommand{\headrulewidth}{0pt}
\setlength{\parindent}{0pt}\setlength{\parskip}{6pt}
\begin{document}
\begin{center}{\LARGE 高一数学拓展练习：教师解析}\end{center}
本卷为按本次教学需求编制的题目，非引用真题；不作题型首创声明。适用前提：已学集合运算、充分必要条件、命题否定及正数基本不等式。难度为相对常规初学练习的设计判断，未经学生实测。共100分，建议60分钟；第5题可作为课后挑战。任何正确的等价解法均应按对应推理给分。

\textbf{1．答案：}\quad（1）$a<\frac32$；（2）$a=0$。

（1）当 $a<0$ 时，$a>2a$，故 $B_a=\varnothing\subseteq A$。
当 $a\ge0$ 时，$B_a=[a,2a]$，包含于 $A=[0,3)$ 当且仅当 $2a<3$。
合并得 $a<\frac32$。

（2）$a<0$ 时交集为空；$a=0$ 时交集为 $\{0\}$。
若 $0<a<3$，因 $a<\min(2a,3)$，在 $a$ 右侧仍有无穷多个元素属于交集；若 $a\ge3$，交集为空。故只有 $a=0$。

评分：第（1）问10分，空集情形4分、非空情形4分、合并2分；第（2）问8分，验证 $a=0$ 得2分，排除 $a<0$、$0<a<3$、$a\ge3$ 各2分。

\textbf{2．答案：}\quad（1）$-2\le a\le-1$；（3）不存在。

（1）设 $P=(-1,2)$，$Q=(a,a+4)$。充分不必要等价于 $P\subsetneq Q$。包含关系要求 $a\le-1$ 且 $a+4\ge2$，即 $-2\le a\le-1$。这些参数下，$Q$ 的区间长度为4，$P$ 的长度为3，两集合不相等，故确为真包含。

（2）存在实数 $x$，使得 $-1<x<2$，且（$x\le a$ 或 $x\ge a+4$）。

（3）若 $Q\subseteq P$，必要且充分的端点条件为 $a\ge-1$ 且 $a+4\le2$，即同时有 $a\ge-1$ 与 $a\le-2$，矛盾。

评分：第（1）问8分，包含方向2分、端点条件及范围4分、排除相等2分；第（2）问6分，“存在”2分、保留 $-1<x<2$ 得2分、结论的否定及“或”2分；第（3）问4分，端点条件2分、矛盾结论2分。

\textbf{3．答案：}\quad（1）最小值为9，$x=\frac13,y=\frac23$；（2）不等式均成立，但最小值结论错误。

（1）利用 $x+y=1$，
\[\frac1x+\frac4y=(x+y)\left(\frac1x+\frac4y\right)=5+\frac yx+\frac{4x}y\ge5+2\sqrt4=9.\]
等号当且仅当 $\frac yx=\frac{4x}y$，由正数条件得 $y=2x$，结合和为1得到 $x=\frac13,y=\frac23$，此时原式确为9。

（2）由基本不等式 $xy\le\frac14$ 可知第二步正确，第一步也正确。但第一步等号要求 $y=4x$，第二步等号要求 $x=y$，在正数条件下不可能同时成立。因此8只是一个下界，并非可以取得的最小值。

评分：第（1）问12分，乘以 $x+y$ 并展开4分、基本不等式4分、等号条件及可达性4分；第（2）问8分，判断两步正确2分、两个等号条件各2分、不能同时成立及结论2分。
\newpage
\textbf{4．答案：}\quad $k\le\sqrt2$；当 $k=\sqrt2$ 时，等号仅在 $x=y=\frac1{\sqrt2}$ 时成立。

必要性：令 $x=y=\frac1{\sqrt2}$，原不等式给出 $2\ge k\sqrt2$，故 $k\le\sqrt2$。

充分性：对任意正数 $x,y$，
\[x^2+\frac12\ge\sqrt2x,\qquad y^2+\frac12\ge\sqrt2y.\]
相加得 $x^2+y^2+1\ge\sqrt2(x+y)$。因 $x+y>0$，当 $k\le\sqrt2$ 时，必有 $\sqrt2(x+y)\ge k(x+y)$。所以全部满足条件的参数为 $(-\infty,\sqrt2]$。

当 $k=\sqrt2$ 时，两处基本不等式同时取等要求 $x^2=y^2=\frac12$，由正数条件得上述唯一数对。

独立复核：$x^2+y^2+1-\sqrt2(x+y)=(x-\frac1{\sqrt2})^2+(y-\frac1{\sqrt2})^2$，直接验证非负性及唯一取等点。

评分：必要性6分；充分性10分，其中证明临界不等式6分、说明所有较小 $k$ 也成立4分；参数范围2分；全部等号条件4分。仅代入一个数对得到上界，未证明对所有正数成立，不能获得充分性分数。

\textbf{5．答案：}\quad（1）$S=\{2,-1,\frac12\}$；（2）不存在。

（1）由 $2\in S$ 得 $-1\in S$，再得 $\frac12\in S$，三个数互异。因恰有三个元素，$S$ 只能是此集合。验证：
\[2\longmapsto-1\longmapsto\frac12\longmapsto2.\]
且集合中没有1，所以满足全部条件。

（2）先说明 $0\notin S$：否则由给定性质会有 $1\in S$，矛盾。
任取 $t\in S$，则 $t\ne0,1$。连续应用给定性质，得到
\[t,\qquad u=\frac1{1-t},\qquad v=\frac{t-1}{t}\quad\text{都属于 }S,\qquad \frac1{1-v}=t.\]
这三个数两两不同。事实上，$t=u$、$u=v$ 或 $v=t$，每一种情况整理后都得到
\[t^2-t+1=0,\]
但 $t^2-t+1=(t-\frac12)^2+\frac34>0$，矛盾。

若 $S$ 恰有四个元素，记第四个为 $s$。由性质知 $\frac1{1-s}\in S$。它不能等于 $s$，因为方程 $s=\frac1{1-s}$ 同样没有实数解；也不能等于 $t,u,v$ 中任何一个：映射 $r\mapsto\frac1{1-r}$ 在 $r\ne1$ 时具有单射性，即输出相等必有输入相等，而 $t,u,v$ 已分别由 $v,t,u$ 映得。于是第四个元素无处映入，矛盾。

评分：第（1）问8分，推出三个元素4分、确定集合2分、验证全部条件2分；第（2）问14分，排除0得2分、推出三步循环3分、证明三个数互异3分、排除第四个数映向自身2分、排除映向前三个数并得结论4分。可用等价的逐一解方程代替“单射性”。

\newpage
\textbf{考查与反馈（教师使用）}

第1题：目标为空集、子集与端点；预期分类讨论，难度中等。常见断点是直接比较端点而遗漏空集，或把3计入交集；依据分类完整性评价。

第2题：目标为条件关系与含量词命题的否定；预期将条件转成集合包含，难度中等。常见错误为包含方向颠倒、端点机械取严格不等号、否定时连前提一起否定；依据包含方向、量词及联结词评价。

第3题：目标为基本不等式与最小值可达性；预期借助定和条件展开，难度中等偏难。常见错误为把任意下界当最小值；依据等号是否可同时实现评价。

第4题：目标为恒成立命题的必要性和充分性；预期临界点检验配合拆分常数，难度较难。常见错误为只检验特殊值；依据是否覆盖任意正数评价。

第5题：目标为读懂新定义并沿规则推理；预期连续代入、配方、反证，难度挑战。无需预先讲授映射术语；依据是否解释元素互异与第四个元素不可能存在评价。可课后提示“从任意元素出发，连续计算三次”。
\end{document}
